\subsection{Inverses}
\label{subsec:inverses}

\begin{definition}[Inverse Element]
\label{def:inverse-element}
Let $(G,\star)$ be a group with identity element $e$.
An element $b\in G$ is an \textbf{inverse} of $a\in G$ if
\[
a\star b=b\star a=e.
\]
\end{definition}

\begin{remark*}[Standard quantified statement]
\[
\operatorname{InverseElement}(b,a,\star,e,G)
\Longleftrightarrow
a,b,e\in G
\land
a\star b=b\star a=e.
\]
\end{remark*}

\begin{remark*}[Definition predicate reading]
$b$ is an inverse of $a$ exactly when composing $a$ and $b$ in either order
returns the identity.
\end{remark*}

\begin{remark*}[Negated quantified statement]
\[
\neg\operatorname{InverseElement}(b,a,\star,e,G)
\Longleftrightarrow
b\notin G
\lor
a\star b\neq e
\lor
b\star a\neq e.
\]
\end{remark*}

\begin{remark*}[Interpretation]
An inverse is an element that undoes another element with respect to the group
operation.
\end{remark*}

\begin{remark*}[Dependencies]
\begin{itemize}
  \item Group.
  \item Identity element.
\end{itemize}
\end{remark*}

\begin{proposition}[Uniqueness of Inverses]
\label{prop:group-inverse-unique}
In a group, every element has a unique inverse.
\end{proposition}

\begin{remark*}[Standard quantified statement]
\[
\operatorname{Group}(G,\star)
\Longrightarrow
\forall a\in G,\; \exists! b\in G\;
\operatorname{InverseElement}(b,a,\star,e,G).
\]
\end{remark*}

\begin{remark*}[Theorem predicate reading]
Every element of a group has exactly one inverse.
\end{remark*}

\begin{remark*}[Negated quantified statement]
\[
\operatorname{Group}(G,\star)
\land
\exists a,b,c\in G,\;
\operatorname{InverseElement}(b,a,\star,e,G)
\land
\operatorname{InverseElement}(c,a,\star,e,G)
\land b\neq c
\Longrightarrow \bot.
\]
\end{remark*}

\begin{remark*}[Interpretation]
Once the group operation and identity are fixed, inverse notation is
well-defined.
\end{remark*}

\begin{remark*}[Dependencies]
\begin{itemize}
  \item Group.
  \item Inverse element.
  \item Identity element.
\end{itemize}
\end{remark*}

\begin{remark*}[Proof TODO]
Proof exists in the archived chapter and should be migrated to a
one-proof-per-file artifact.
Source: \texttt{archive/algebra-and-abstract-structures-legacy.tex}.
\end{remark*}
